\documentclass{article}

\usepackage{amsfonts}
\usepackage{amsthm}
\usepackage{mathrsfs}
\usepackage{amsmath}
\usepackage{graphicx}
\usepackage{hyperref}
\usepackage{stackengine}
\usepackage[mathscr]{eucal}
\usepackage{listings}
\usepackage{algorithm}
\usepackage{algpseudocode}
\usepackage{float}
\usepackage{xcolor}
\usepackage{url}
\usepackage{amssymb}
\newtheorem{theorem}{Theorem}
\theoremstyle{definition}\newtheorem{definition}{DEF}
\theoremstyle{lemma}\newtheorem{lemma}{Lemma}
\theoremstyle{corrolary}\newtheorem{corrolary}{Cor}
\setlength{\parindent}{0pt}
\definecolor{codebg}{RGB}{248,248,248}
\definecolor{codekeyword}{RGB}{0,76,153}
\definecolor{codecomment}{RGB}{0,128,96}
\definecolor{codestring}{RGB}{163,21,21}
\definecolor{codenumber}{RGB}{102,102,102}
\lstdefinestyle{pythoncode}{
    language=Python,
    basicstyle=\ttfamily\small,
    keywordstyle=\color{codekeyword}\bfseries,
    commentstyle=\color{codecomment}\itshape,
    stringstyle=\color{codestring},
    numberstyle=\scriptsize\color{codenumber},
    numbers=left,
    numbersep=10pt,
    stepnumber=1,
    showstringspaces=false,
    breaklines=true,
    keepspaces=true,
    columns=fullflexible,
    tabsize=4,
    frame=single,
    framerule=0pt,
    rulecolor=\color{codebg},
    backgroundcolor=\color{codebg}
}
\lstset{style=pythoncode}

\begin{document}

\begin{definition}
A complex random variable is a measurable function $Z: (S, \Omega, p) \rightarrow \mathbb{C}$.
\end{definition}

\begin{lemma}
Let $Z$ be a complex random variable.
Then $\mathcal{R}(Z)$ and $\mathcal{I}(Z)$ are random varaibles.
\end{lemma}
\begin{proof}
Notice $\mathbb{R}\subseteq \mathcal{B}(\mathbb{C})$.
So for any $B \in \mathcal{B}(\mathbb{C})$, $B\cap \mathbb{R}\in \mathcal{B}(\mathbb{C})$.
So $(\mathcal{R}\circ Z)^{-1} (B) = Z^{-1}(B\cap \mathbb{R}) \in \Omega$.
Thus, $\mathcal{R} \circ Z: (S, \Omega , p) \rightarrow \mathbb{R}$ is a measurable function.
So it is  a ranom vairable.
Simiilar logic applies to $\mathcal{I}\circ Z$.
\end{proof}

\begin{definition}
Let $Z=XX+iY$ be a complex random variable.
Say the complex CDF of $Z$ is
\begin{equation}
F(x+iy)=\Pr(X\leq x, Y \leq y)
\end{equation}
Say the complex PDF is
\begin{equation}
f(x+iy)=\frac{\partial^2 F}{\partial x \partial y}
\end{equation}
\end{definition}

\begin{lemma}
Let $Z=X+iY$ be a complex random variable such that the expected value exists.
Then the expected value is
\begin{equation}
\mathbb{E}(Z) = \mathbb{E}(X) + i\mathbb{E}(Y)
\end{equation}
\end{lemma}
\begin{proof}
This is a consequence of the linearity of expectation.
\end{proof}

\begin{lemma}
\begin{equation}
\mathbb{E}(\bar{Z}) = \overline{\mathbb{E}(Z)}
\end{equation}
\end{lemma}
\begin{proof}
Let $X=\mathcal{R}(Z)$ and $Y=\mathcal{I}(Z)$.
Then $\bar{Z}=X-iY$.
So
\begin{align*}
\mathbb{E}(\bar{Z})
&= \mathbb{E}(X-iY)\\
&=\mathbb{E}(X) -i\mathbb{E}(Y)\\
&=\overline{\mathbb{E}(Z)}
\end{align*}
\end{proof}

\begin{definition}
Let $Z$ be a complex random variable with expected value $\mu$.
Say the variance of a complex random variable is
\begin{equation}
Var(Z) = \mathbb{E}(\overline{(Z-\mu)}(Z-\mu))
\end{equation}
\end{definition}

\begin{lemma}
\begin{equation}
Var(Z) = \mathbb{E}(|Z|^2) - (\mathbb{E}(|Z|))^2
\end{equation}
\end{lemma}
\begin{proof}
\begin{align*}
Var(Z)
&= \mathbb{E}(\overline{(Z - \mu)}(Z - \mu)) \\
&= \mathbb{E}((\bar{Z} - \bar{\mu})(Z-\mu)) \\
&=\mathbb{E}(\bar{Z}Z -\bar{Z}\mu -\bar{\mu}Z + \bar{\mu}\mu) \\
&=\mathbb{E}(\bar{Z}Z) -\overline{\mathbb{E}(Z)}\mu - \bar{\mu}\mathbb{E}(Z) + \bar{\mu}\mu  \\
&= \mathbb{E}(|Z|^2) - |\mu|^2
\end{align*}
\end{proof}

\begin{definition}
Let $Z$ be a complex random variable with mean $\mu$.
Say teh pseudo-variance is
\begin{equation}
pVar(Z) = \mathbb{E}((Z-\mu)^2)
\end{equation}
\end{definition}

\begin{lemma}
\begin{equation}
pVar(Z) = \mathbb{E}(Z^2) - (\mathbb{E}(Z))^2
\end{equation}
\end{lemma}
\begin{proof}
\begin{align*}
pVar(Z)
&= \mathbb{E}((Z-\mu)^2) \\
&= \mathbb{E}(Z^2 - 2\mu Z + \mu^2) \\
&= \mathbb{E}(Z^2) - \mu^2
\end{align*}
\end{proof}

\textbf{Note-}
The variance of a complex random variable will always be a postivie real number but the pseudo-variance can be any complex number.

\begin{definition}
Say a measurable function $Z: (S, \Omega, p)\rightarrow \mathbb{C}^n$ is a complex random vector.
\end{definition}

\begin{definition}
The covariance of two complex randomm variables is
\begin{equation}
Cov(Z, W) = \mathbb{E}(\overline{(Z - \mathbb{E}(Z))}(W-\mathbb{E}(W)))
\end{equation}
\end{definition}

\begin{lemma}
\begin{equation}
Cov(Z, Z) = Var(Z)
\end{equation}
\end{lemma}
\begin{proof}
This follows from the definitions of variance and covariance.
\end{proof}

\begin{lemma}
\begin{equation}
Cov(Z, W) = \mathbb{E}(\bar{Z}W) - \mathbb{E}(\bar{Z}) \mathbb{E}(W)
\end{equation}
\end{lemma}
\begin{proof}
Notice that for any random variables $W$ and $Z$,
\begin{align*}
\mathbb{E}(\overline{(Z - \mathbb{E}(Z))}(W - \mathbb{E}(W)))
&= \mathbb{E}(\bar{Z}W -\mathbb{E}(\bar{Z})W -\bar{Z}\mathbb{E}(W) + \mathbb{E}(\bar{Z})\mathbb{E}(W)) \\
&= \mathbb{E}(\bar{Z}W) -2\mathbb{E}(\bar{Z})\mathbb{E}(W) + \mathbb{E}(\bar{Z})\mathbb{E}(W)\\
&= \mathbb{E}(\bar{Z}W)-\mathbb{E}(\bar{Z})\mathbb{E}(W)
\end{align*}
\end{proof}

\begin{theorem}
\begin{equation}
Cov(Z,W)=\overline{Cov(W,Z)}
\end{equation}
\end{theorem}
\begin{proof}
\begin{align*}
Cov(Z,W)
&= \mathbb{E}(\overline{(Z - \mathbb{E}(Z))}(W - \mathbb{E}(W))) \\
&= \overline{\mathbb{E}((Z - \mathbb{E}(Z))\overline{(W - \mathbb{E}(W))})} \\
&= \overline{\mathbb{E}(\overline{(W - \mathbb{E}(W))}(Z - \mathbb{E}(Z)))} \\
&=\overline{Cov(W,Z)}
\end{align*}
\end{proof}

\begin{theorem}
\begin{equation}
Var(Z + W) = Var(Z) + Var(W) + Cov(Z,W) + Cov(W,Z)
\end{equation}
\end{theorem}
\begin{proof}
\begin{align*}
Var(Z+W)
&= \overline{(W + Z - \mathbb{E}(W + Z))}(W + Z - \mathbb{E}(W + Z)) \\
&=\overline{(Z - \mathbb{E}(Z) + W - \mathbb{E}(W))}(Z - \mathbb{E}(Z) + W - \mathbb{E}(W)) \\
&= (\overline{(Z - \mathbb{E}(Z))}(Z - \mathbb{E}(Z)) + \overline{Z - \mathbb{E}(Z))}(W - \mathbb{E}(W)) \\
&\quad \quad \quad \quad + \overline{(W - \mathbb{E}(W))}(Z - \mathbb{E}(Z)) \\
&\quad \quad \quad \quad + \overline{(W - \mathbb{E}(W))}(W - \mathbb{E}(W))
\end{align*}
Taking the expected value gives the desired result.
\end{proof}

\begin{theorem}
Let $Z_1, Z_2, W$ be complex random variables and $a,b\in \mathbb{C}$.
\begin{equation}
Cov(W, aZ_1 + bZ_2) = a Cov(W, Z_1) + b\, Cov(W, Z_2)
\end{equation}
\end{theorem}
\begin{proof}
By linearity of expectation, $\mathbb{E}(aZ_1 + bZ_2) = a\mathbb{E}(Z_1) + b\mathbb{E}(Z_2)$.
Thus,
\begin{equation*}
\mathbb{E}(\bar{W})\mathbb{E}(aZ_1 + bZ_2) = a\mathbb{E}(\bar{W})\mathbb{E}(Z_1) + b\mathbb{E}(\bar{W})\mathbb{E}(Z_2)
\end{equation*}
Similarly,
\begin{equation*}
\mathbb{E}(\bar{W}(aZ_1 + bZ+_2)) = a\mathbb{E}(\bar{W}Z_1) + b\mathbb{E}(\bar{W}Z_2)
\end{equation*}
Thus,
\begin{align*}
Cov(W, aZ_1 + bZ_2)
&=\mathbb{E}(\bar{W}(aZ_1 + bZ+_2)) - \mathbb{E}(\bar{W})\mathbb{E}(aZ_1 + bZ_2) \\
&= a\mathbb{E}(\bar{W}Z_1) + b\mathbb{E}(\bar{W}Z_2) -a\mathbb{E}(\bar{W})\mathbb{E}(Z_1) - b\mathbb{E}(\bar{W})\mathbb{E}(Z_2)
&=aCov(W,Z_1) + b\, Cov(W, Z_2)
\end{align*}
\end{proof}
\begin{corrolary}
\begin{equation}
Cov(aZ_1 + bZ_2, W) = \bar{a}Cov(Z_1, W) + \bar{b}\, Cov(Z_2, W)
\end{equation}
\end{corrolary}

\begin{definition}
The pseudocovariance of two complex random variables is
\begin{equation}
pCov(Z,W)=\mathbb{E}(Z-\mathbb{E}(Z))(W-\mathbb{E}(W)))
\end{equation}
\end{definition}

\begin{lemma}
\begin{equation}
pCov(Z,W) = Cov(\bar{Z},W)
\end{equation}
\end{lemma}
\begin{proof}
\begin{align*}
pCov(Z,W)
&= \mathbb{E}((Z -\mathbb{E}(Z))(W - \mathbb{E}(W)))\\
&= \mathbb{E}(\overline{(\bar{Z} - \mathbb{E}(\bar{Z}))}(W - \mathbb{E}(W))) \\
&= Cov(\bar{Z}, W)
\end{align*}
\end{proof}
\begin{corrolary}
\begin{equation}
pCov(Z,W) = \mathbb{E}(ZW) - \mathbb{E}(Z)\mathbb{E}(W)
\end{equation}
\end{corrolary}
\begin{corrolary}
\begin{equation}
pCov(Z, Z) = pVar(Z)
\end{equation}
\end{corrolary}
\begin{corrolary}
\begin{equation}
pVar(Z+W) = pVar(Z) + pVar(W) + 2pCov(Z,W)
\end{equation}
\end{corrolary}
\begin{corrolary}
\begin{equation}
pCov(W, aZ_1 + bZ_2) = apCov(W,Z_1) + b\, pCov(W,Z_2)
\end{equation}
\end{corrolary}

\begin{definition}
Say a complex random variable $Z$ is proper whenever $\mathbb{E}(Z)=0$, $\mathbb{E}(Z^2)=0$, and $Var(Z)<\infty$.
\end{definition}

\textbf{Note-}
Proper random variables are a generalization ofrandom variables with finite variance.
The requirement that $\mathbb{E}(Z^2)=0$ preserves some important features that will not generalize to the complex plane with only 0 mean and finite variance.

\begin{lemma}
If $Z=X+iY$ is a proper complex random variable then $Cov(X,Y)=0$.
\end{lemma}
\begin{proof}
Notice $Z^2 = (X+iY)^2 = (X^2 - Y^2) + 2iXY$.
Let $Z$ be proper, so $\mathbb{E}(Z^2)=0$.
Then the real and imaginary parts must both also be 0, so $\mathbb{E}(XY)=0$.
Also, $Z$ has 0 mean, so $X$ and $Y$ are also both 0 mean.
Thus, $Cov(X,Y)=0$.
\end{proof}

\begin{definition}
Say a complex random variabel $Z$ is circularly symetric whenever $e^{i\theta}Z$ has the same distribution for all $\theta$.
\end{definition}

\begin{theorem}
If $Z$ is a circularly symmetric random variable with finite variance then it is proper.
\end{theorem}
\begin{proof}
Let $Z$ be a circularly symmetric random variable with $Var(Z)<\infty$.
Then $e^{i\theta}Z$ has the same distribution for all $\theta$.
In particular, consider $\theta=\frac{\pi}{2}$.
Then $\mathbb{E}(Z)=\mathbb{E}(iZ)=i\mathbb{E}(Z)$.
Thus, $\mathbb{E}(Z)=0$.
Similarly, $\mathbb{E}(Z^2)=\mathbb{E}(-Z^2)=-\mathbb{E}(Z^2)$.
So $\mathbb{E}(Z^2)=0$.
Therefore, $Z$ must be proper.
\end{proof}

\begin{definition}
Say the covariance matrix of a complex random vector $Z$ is
\begin{equation}
Cov(Z) = \mathbb{E}((Z - \mathbb{E}(Z))(Z - \mathbb{E}(Z))^H)
\end{equation}
\end{definition}

\begin{definition}
Say the pseudo-covariance matrix of a complex random vector is
\begin{equation}
pCov(Z) = \mathbb{E}((Z - \mathbb{E}(Z))(Z - \mathbb{E}(Z))^T)
\end{equation}
\end{definition}

\begin{lemma}
The diagonal elements of $Cov(Z)$ are $Var(Z_i)$ and the off diagonal ellements are $Cov(Z_i, Z_j)$.
The diagonal elements of $pCov(Z)$ are $pVar(Z_i)$ and the off diagonal elements are $pCov(Z_i, Z_j)$.
\end{lemma}
\begin{proof}
This follows from the definition of teh covariance and pseudo-covariance matrices.
\end{proof}

\end{document}
